% Paso 7: Plantilla de Síntesis Interpretativa
% Proyecto: Análisis Exegético Profundo
% Ministerio "La Gloria es del Señor"

\documentclass[11pt,a4paper]{article}
\usepackage[utf8]{inputenc}
\usepackage[spanish]{babel}
\usepackage[margin=2.5cm]{geometry}
\usepackage{graphicx}
\usepackage{fancyhdr}
\usepackage{xcolor}
\usepackage{tcolorbox}
\usepackage{enumitem}
\usepackage{tabularx}
\usepackage{multirow}
\usepackage{amsfonts}
\usepackage{fontspec}
\usepackage{titlesec}

% Configuración de colores
\definecolor{forestgreen}{RGB}{27,67,50}
\definecolor{sagegreen}{RGB}{45,90,39}
\definecolor{olivegreen}{RGB}{64,145,108}
\definecolor{mintgreen}{RGB}{82,183,136}
\definecolor{lightmint}{RGB}{116,198,157}
\definecolor{softcream}{RGB}{216,243,220}
\definecolor{papergray}{RGB}{118,128,150}

% Corrección para el warning de fancyhdr
\setlength{\headheight}{14pt}

% Configuración de encabezados
\pagestyle{fancy}
\fancyhf{}
\fancyhead[L]{\textcolor{forestgreen}{\textbf{Análisis Exegético Profundo}}}
\fancyhead[R]{\textcolor{papergray}{Paso 7: Síntesis Interpretativa}}
\fancyfoot[C]{\textcolor{papergray}{\thepage}}
\fancyfoot[R]{\textcolor{papergray}{\scriptsize Ministerio "La Gloria es del Señor"}}

% Configuración de títulos
\titleformat{\section}
{\color{forestgreen}\Large\bfseries}
{\thesection}{1em}{}

\titleformat{\subsection}
{\color{sagegreen}\large\bfseries}
{\thesubsection}{1em}{}

% Configuración de cajas
\tcbuselibrary{skins}
\newtcolorbox{infobox}[1]{
    colback=softcream,
    colframe=olivegreen,
    fonttitle=\bfseries,
    title=#1,
    left=10pt,
    right=10pt,
    top=10pt,
    bottom=10pt
}

\newtcolorbox{checklistbox}{
    colback=white,
    colframe=mintgreen,
    left=15pt,
    right=15pt,
    top=10pt,
    bottom=10pt,
    boxrule=2pt
}

\newtcolorbox{synthesisbox}{
    colback=lightmint!20,
    colframe=sagegreen,
    left=10pt,
    right=10pt,
    top=8pt,
    bottom=8pt,
    boxrule=1.5pt
}

\newtcolorbox{interpretationbox}{
    colback=olivegreen!10,
    colframe=olivegreen,
    left=10pt,
    right=10pt,
    top=8pt,
    bottom=8pt,
    boxrule=1.5pt
}

\newtcolorbox{theologicalbox}{
    colback=forestgreen!10,
    colframe=forestgreen,
    left=10pt,
    right=10pt,
    top=8pt,
    bottom=8pt,
    boxrule=1.5pt
}

\begin{document}

% Título principal
\begin{center}
    {\Huge\color{forestgreen}\textbf{PASO 7}}\\[0.5cm]
    {\Large\color{sagegreen}\textbf{SÍNTESIS INTERPRETATIVA}}\\[0.3cm]
    {\large\color{papergray}Integración de Hallazgos en Interpretación Coherente}\\[1cm]
    
    \begin{tcolorbox}[colback=olivegreen!10, colframe=olivegreen, width=0.8\textwidth]
        \centering
        \textcolor{forestgreen}{\textbf{Proyecto: Análisis Exegético Profundo}}\\
        \textcolor{papergray}{Metodología de 8 Pasos para Estudio Bíblico Riguroso}
    \end{tcolorbox}
\end{center}

\vspace{1cm}

% Información del estudiante
\begin{infobox}{Información del Proyecto}
\begin{tabularx}{\textwidth}{|X|X|}
\hline
\textbf{Nombre del Estudiante:} & \rule{6cm}{0.4pt} \\[0.8cm]
\hline
\textbf{Texto Bíblico:} & \rule{6cm}{0.4pt} \\[0.8cm]
\hline
\textbf{Fecha del Análisis:} & \rule{6cm}{0.4pt} \\[0.8cm]
\hline
\textbf{Horas de Investigación:} & \rule{6cm}{0.4pt} \\[0.8cm]
\hline
\end{tabularx}
\end{infobox}

\section{Objetivo del Paso 7: Síntesis Interpretativa}

La síntesis interpretativa integra todos los datos exegéticos recolectados en los pasos anteriores para formular una interpretación coherente y bien fundamentada. Este es el momento de unir observación, análisis léxico, gramática, contexto histórico, género literario y contexto canónico en una comprensión unificada del texto.

\subsection{Descripción del Proceso}

Sintetizaremos todos los hallazgos previos, identificaremos el mensaje principal del texto, resolveremos cualquier tensión interpretativa, y formularemos la interpretación más probable. También consultaremos comentarios de referencia para validar y refinar nuestra comprensión.

\begin{infobox}{Principio Hermenéutico Central}
La interpretación debe estar basada en evidencia textual sólida y ser consistente con todos los datos exegéticos. Busca la explicación más simple que mejor explique toda la evidencia.
\end{infobox}

\section{Lista de Verificación}

\begin{checklistbox}
\textbf{Marca cada elemento completado:}

\begin{itemize}[leftmargin=1cm]
    \item[$\square$] Sintetizar todos los hallazgos de los pasos 1-6
    \item[$\square$] Formular el mensaje principal del texto
    \item[$\square$] Identificar puntos teológicos clave
    \item[$\square$] Resolver tensiones interpretativas aparentes
    \item[$\square$] Verificar coherencia interna de la interpretación
    \item[$\square$] Consultar comentarios académicos de referencia
    \item[$\square$] Evaluar interpretaciones alternativas
    \item[$\square$] Formular declaración interpretativa final
\end{itemize}
\end{checklistbox}

\newpage

\section{Resumen de Hallazgos Previos}

\subsection{Síntesis de los Pasos 1-6}

\begin{synthesisbox}
\textbf{Integra aquí los descubrimientos más importantes de cada paso:}
\end{synthesisbox}

\textbf{Paso 1 - Observación Inicial:}
\textit{(Elementos clave observados, preguntas principales, estructura)}
\vspace{2cm}

\textbf{Paso 2 - Análisis Léxico:}
\textit{(Palabras clave más significativas, matices importantes)}
\vspace{2cm}

\textbf{Paso 3 - Análisis Gramatical:}
\textit{(Estructuras gramaticales clave, flujo del argumento)}
\vspace{2cm}

\textbf{Paso 4 - Contexto Histórico-Cultural:}
\textit{(Trasfondo más relevante, factores culturales importantes)}
\vspace{2cm}

\textbf{Paso 5 - Género Literario:}
\textit{(Género identificado, características que afectan interpretación)}
\vspace{2cm}

\textbf{Paso 6 - Contexto Canónico:}
\textit{(Relación con otros textos, lugar en la revelación bíblica)}
\vspace{2cm}

\newpage

\section{Identificación del Mensaje Central}

\subsection{Idea Principal del Texto}

\begin{interpretationbox}
\textbf{En una oración clara, ¿cuál es el mensaje principal que el autor quería comunicar a su audiencia original?}
\end{interpretationbox}

\vspace{4cm}

\subsection{Temas Secundarios}

\textbf{¿Qué temas o ideas secundarias apoyan el mensaje principal?}

\begin{enumerate}[leftmargin=1cm]
    \item \rule{12cm}{0.4pt}
    \vspace{0.8cm}
    \item \rule{12cm}{0.4pt}
    \vspace{0.8cm}
    \item \rule{12cm}{0.4pt}
    \vspace{0.8cm}
    \item \rule{12cm}{0.4pt}
    \vspace{0.8cm}
\end{enumerate}

\subsection{Propósito del Autor}

\textbf{¿Con qué propósito específico escribió el autor este texto?}
\begin{itemize}[leftmargin=1cm]
    \item[$\square$] Enseñar/Instruir
    \item[$\square$] Corregir/Disciplinar
    \item[$\square$] Consolar/Animar
    \item[$\square$] Advertir/Prevenir
    \item[$\square$] Motivar/Exhortar
    \item[$\square$] Explicar/Clarificar
    \item[$\square$] Defender/Argumentar
    \item[$\square$] Otro: \rule{6cm}{0.4pt}
\end{itemize}

\textbf{Explica cómo el texto cumple este propósito:}
\vspace{3cm}

\section{Análisis Teológico}

\subsection{Doctrinas Principales}

\begin{theologicalbox}
\textbf{¿Qué doctrinas bíblicas fundamentales toca o desarrolla este texto?}
\end{theologicalbox}

\begin{tabularx}{\textwidth}{|X|X|}
\hline
\textbf{Doctrina} & \textbf{Cómo la Desarrolla el Texto} \\
\hline
Teología Propia (Dios) & \\[1.5cm]
\hline
Cristología (Cristo) & \\[1.5cm]
\hline
Pneumatología (Espíritu Santo) & \\[1.5cm]
\hline
Soteriología (Salvación) & \\[1.5cm]
\hline
Eclesiología (Iglesia) & \\[1.5cm]
\hline
Escatología (Últimas Cosas) & \\[1.5cm]
\hline
Otras: & \\[1.5cm]
\hline
\end{tabularx}

\subsection{Contribución Teológica Única}

\textbf{¿Qué perspectiva única o énfasis especial aporta este texto a estas doctrinas?}
\vspace{3cm}

\textbf{¿Hay algún aspecto teológico que este texto clarifica o desarrolla de manera distintiva?}
\vspace{3cm}

\newpage

\section{Resolución de Tensiones Interpretativas}

\subsection{Dificultades Identificadas}

\textbf{¿Qué aspectos del texto presentaron dificultades interpretativas?}

\begin{tabularx}{\textwidth}{|X|X|X|}
\hline
\textbf{Dificultad} & \textbf{Opciones Interpretativas} & \textbf{Resolución Adoptada} \\
\hline
& & \\[2cm]
\hline
& & \\[2cm]
\hline
& & \\[2cm]
\hline
\end{tabularx}

\subsection{Justificación de Decisiones Interpretativas}

\textbf{¿Por qué elegiste estas interpretaciones sobre las alternativas?}
\vspace{3cm}

\textbf{¿Qué evidencia textual y contextual apoya tus decisiones?}
\vspace{3cm}

\section{Verificación de Coherencia}

\subsection{Consistencia Interna}

\textbf{¿Tu interpretación explica satisfactoriamente todos los elementos del texto?}
\vspace{2cm}

\textbf{¿Hay algún aspecto del texto que tu interpretación no aborde adecuadamente?}
\vspace{2cm}

\textbf{¿Tu interpretación es consistente con el género literario del texto?}
\vspace{2cm}

\subsection{Consistencia Contextual}

\textbf{¿Tu interpretación armoniza con el contexto inmediato?}
\vspace{2cm}

\textbf{¿Es consistente con el mensaje del libro completo?}
\vspace{2cm}

\textbf{¿Está en línea con la teología del autor en otros escritos?}
\vspace{2cm}

\textbf{¿Armoniza con el testimonio general de la Escritura?}
\vspace{2cm}

\newpage

\section{Consulta de Comentarios}

\subsection{Comentarios Consultados}

\begin{tabularx}{\textwidth}{|X|X|X|}
\hline
\textbf{Comentario/Autor} & \textbf{Interpretación Principal} & \textbf{Puntos de Acuerdo/Diferencia} \\
\hline
& & \\[2cm]
\hline
& & \\[2cm]
\hline
& & \\[2cm]
\hline
& & \\[2cm]
\hline
\end{tabularx}

\subsection{Evaluación de Interpretaciones Alternativas}

\textbf{¿Qué interpretaciones alternativas consideraste y por qué las rechazaste?}
\vspace{3cm}

\textbf{¿En qué puntos tus comentarios favoritos están en desacuerdo, y cómo resolviste estas diferencias?}
\vspace{3cm}

\section{Declaración Interpretativa Final}

\subsection{Mensaje Principal (Reformulado)}

\begin{interpretationbox}
\textbf{Basado en todo tu análisis, reformula el mensaje principal del texto en 2-3 oraciones claras:}
\end{interpretationbox}

\vspace{4cm}

\subsection{Puntos Teológicos Clave}

\textbf{Resume los 3 puntos teológicos más importantes que emergen del texto:}

\begin{enumerate}[leftmargin=1cm]
    \item \textbf{Punto Teológico \#1:}
    \vspace{2cm}
    
    \item \textbf{Punto Teológico \#2:}
    \vspace{2cm}
    
    \item \textbf{Punto Teológico \#3:}
    \vspace{2cm}
\end{enumerate}

\subsection{Significado para la Audiencia Original}

\textbf{¿Qué habría significado este mensaje para la audiencia original?}
\vspace{3cm}

\textbf{¿Cómo habrían respondido a este mensaje en su contexto histórico?}
\vspace{3cm}

\newpage

\section{Preparación para la Aplicación}

\subsection{Principios Universales}

\textbf{¿Qué principios universales y transculturales emergen del texto?}

\begin{enumerate}[leftmargin=1cm]
    \item \rule{12cm}{0.4pt}
    \vspace{0.8cm}
    \item \rule{12cm}{0.4pt}
    \vspace{0.8cm}
    \item \rule{12cm}{0.4pt}
    \vspace{0.8cm}
    \item \rule{12cm}{0.4pt}
    \vspace{0.8cm}
\end{enumerate}

\subsection{Relevancia Contemporánea}

\textbf{¿Qué aspectos del mensaje original son especialmente relevantes para la iglesia contemporánea?}
\vspace{3cm}

\textbf{¿Hay algún aspecto que era específico para la situación original y que no se aplica directamente hoy?}
\vspace{3cm}

\section{Reflexión sobre el Proceso}

\subsection{Crecimiento en la Comprensión}

\textbf{¿Cómo cambió tu comprensión del texto a través del proceso exegético?}
\vspace{3cm}

\textbf{¿Qué aspecto del análisis te ayudó más a entender el texto?}
\vspace{2cm}

\textbf{¿Qué preguntas iniciales fueron respondidas y cuáles persisten?}
\vspace{3cm}

\begin{infobox}{Recordatorio Metodológico}
Una buena interpretación debe estar basada en evidencia sólida, ser internamente coherente, y estar en armonía con el contexto y la teología bíblica general. La humildad intelectual reconoce los límites de nuestro entendimiento.
\end{infobox}

\section{Siguiente Paso}

Con la síntesis interpretativa completada, procederás al \textbf{Paso 8: Aplicación Contemporánea}, donde desarrollaremos aplicaciones relevantes y legítimas para la iglesia de hoy.

\vspace{1cm}

\begin{center}
\begin{tcolorbox}[colback=mintgreen!20, colframe=mintgreen, width=0.7\textwidth]
\centering
\textcolor{forestgreen}{\textbf{¡Interpretación Sólida Lograda!}}\\
\textcolor{papergray}{Has integrado exitosamente todos los datos exegéticos en una comprensión coherente.}
\end{tcolorbox}
\end{center}

\end{document}